\section{Смешанные стратегии}

В классической теории игр предполагается, что игрок может использовать генератор случайных чисел для выбора одной из своих чистых стратегий. Распределение вероятностей, заданное на множестве чистых стратегий, называется смешанной стратегией. Предполагается, что другим игрокам известно данное распределение вероятностей, но не конкретная реализованная чистая стратегия. Заметим, что чистая стратегия является частным случаем смешанной стратегии, при котором один из элементов множества стратегий играется с вероятностью, равной единице.

Для формализации понятия смешанной стратегии вводится понятие симплекс. Пусть $N$ — натуральное число. Тогда $(N-1)$-мерный симплекс описывает все возможные распределения вероятностей на множестве из $N$ элементов и определяется как множество $\Delta^{N-1} \subset \mathbb{R}^N$, состоящее из всех векторов $(p_1, \dots, p_N) \in \Delta^{N-1}$, удовлетворяющих условиям:
\begin{equation}
    \sum_{j=1}^N p_j = 1, \quad p_j \ge 0.
\end{equation}

Рассмотрим конечную игру в нормальной форме $G = \langle I, S, u \rangle$. Назовем $\Delta^{|S_i|-1}$ множеством смешанных стратегий игрока $i$, а декартово произведение $\prod_{i=1}^N \Delta^{|S_i|-1}$ — множеством профилей смешанных стратегий. Элементы $\sigma_i \in \Delta^{|S_i|-1}$ и $\sigma \in \prod_{i=1}^N \Delta^{|S_i|-1}$ называются смешанными стратегиями и профилями смешанных стратегий соответственно. 

При анализе игр в смешанных стратегиях вводятся следующие предположения:
\begin{itemize}
    \item Игрок имеет возможность предоставить выбор чистой стратегии воле случая, но при этом жестко контролировать вероятность реализации той или иной чистой стратегии.
    \item Действия любых двух игроков являются стохастически независимыми событиями.
\end{itemize}

Вероятность того, что при профиле $\sigma$ будет сыгран конкретный профиль чистых стратегий $s$, вычисляется как произведение вероятностей:
\begin{equation}
    \sigma(s) = \prod_{i=1}^N \sigma_i(s_i)
\end{equation}

Тогда математическое ожидание выигрыша игрока $i$ при профиле смешанных стратегий $\sigma$ задается формулой:
\begin{equation}
    u_i(\sigma) = u_i(\sigma_i, \sigma_{-i}) = \sum_{s \in S} \sigma(s)u_i(s)
\end{equation}

\begin{definition}[Смешанное расширение игры]
Пусть $G = \langle I, S, u \rangle$ — конечная игра в нормальной форме. Назовем игру $\langle I, \prod_{i=1}^N \Delta^{|S_i|-1}, u \rangle$ смешанным расширением игры $G$. Тогда равновесие в смешанных стратегиях в игре $G$ — это равновесие Нэша в ее смешанном расширении, где функция $u \colon \prod_{i=1}^N \Delta^{|S_i|-1} \to \mathbb{R}$ является естественным линейным продолжением функции $u \colon S \to \mathbb{R}$.
\end{definition}

\subsection{Существование равновесия}

\begin{theorem}[Дж. Нэш, 1950]\label{thm:nash_existence}
Пусть $G$ — конечная игра. Тогда в игре $G$ существует как минимум одно равновесие в смешанных стратегиях.
\end{theorem}

Для доказательства данной теоремы применяется аппарат теории неподвижных точек, в частности теорема Какутани. Заметим, что если $S$ — некоторое множество, то через $2^S$ обозначается множество всех его подмножеств (булеан).

\begin{theorem}[Какутани, 1941]
Пусть $S$ — непустое, компактное и выпуклое подмножество евклидова пространства. Пусть $r \colon S \to 2^S$ — точечно-множественное отображение, удовлетворяющее следующим условиям:
\begin{enumerate}
    \item Значение $r(s)$ является непустым, выпуклым и компактным множеством для всех $s \in S$.
    \item Отображение $r(\cdot)$ имеет замкнутый график. То есть для любой сходящейся последовательности $(s^n, \hat{s}^n) \to (s, \hat{s})$, такой что $\hat{s}^n \in r(s^n)$ для всех $n$, выполняется $\hat{s} \in r(s)$.
\end{enumerate}
Тогда $r(\cdot)$ имеет неподвижную точку, то есть существует $s^* \in S$, такое что $s^* \in r(s^*)$.
\end{theorem}

\begin{proof}[Доказательство Теоремы \ref{thm:nash_existence}]
Определим функцию реакции (наилучшего ответа) игрока $i$ для смешанных стратегий:
\begin{equation}
    \check{\sigma}_i(\sigma) = \operatorname{argmax}_{\sigma'_i \in \Delta^{|S_i|-1}} u_i(\sigma'_i, \sigma_{-i})
\end{equation}
Данное точечно-множественное отображение определяет смешанные стратегии, максимизирующие ожидаемый выигрыш игрока $i$ при фиксированном профиле стратегий остальных игроков. 
Определим совместное отображение:
\begin{equation}
    \check{\sigma}(\sigma) = \prod_{i=1}^N \check{\sigma}_i(\sigma)
\end{equation}
Заметим, что если $\sigma^*$ является неподвижной точкой отображения $\check{\sigma}(\cdot)$, то $\sigma^*$ по определению является равновесием Нэша. Проверим выполнение условий теоремы Какутани.

\textit{1. Непустота и компактность.} Множество профилей смешанных стратегий выпукло и компактно, так как является декартовым произведением симплексов. Функция ожидаемого выигрыша $u_i(\sigma)$ линейна по $\sigma_i$, следовательно, непрерывна. По теореме Вейерштрасса, непрерывная функция достигает своего максимума на компактном множестве. Значит, $\check{\sigma}_i(\sigma) \neq \emptyset$ для всех профилей $\sigma$.

\textit{2. Выпуклость отображения.} Докажем, что $\check{\sigma}(\sigma)$ выпукло. Пусть $\sigma', \sigma'' \in \check{\sigma}(\sigma)$. Тогда для любого игрока $i$ и любого $\lambda \in (0, 1)$ в силу линейности математического ожидания получаем:
\begin{equation}
    u_i(\lambda\sigma'_i + (1 - \lambda)\sigma''_i, \sigma_{-i}) = \lambda u_i(\sigma'_i, \sigma_{-i}) + (1 - \lambda)u_i(\sigma''_i, \sigma_{-i})
\end{equation}
Поскольку $\sigma'$ и $\sigma''$ являются наилучшими ответами, $u_i(\sigma'_i, \sigma_{-i}) = u_i(\sigma''_i, \sigma_{-i})$. Следовательно:
\begin{equation}
    u_i(\lambda\sigma'_i + (1 - \lambda)\sigma''_i, \sigma_{-i}) = u_i(\sigma'_i, \sigma_{-i})
\end{equation}
Это означает, что выпуклая комбинация $\lambda\sigma'_i + (1 - \lambda)\sigma''_i \in \check{\sigma}_i(\sigma)$, и множество $\check{\sigma}(\sigma)$ выпукло.

\textit{3. Замкнутость графика.} Допустим противное: отображение $\check{\sigma}(\cdot)$ не имеет замкнутого графика. Тогда существует последовательность $(\sigma^n, \hat{\sigma}^n) \to (\sigma, \hat{\sigma})$, такая что $\hat{\sigma}^n \in \check{\sigma}(\sigma^n)$ для всех $n$, но $\hat{\sigma} \notin \check{\sigma}(\sigma)$.
Следовательно, для некоторого игрока $i$ выполняется $\hat{\sigma}_i \notin \check{\sigma}_i(\sigma)$.
Тогда существует альтернативная стратегия $\sigma'_i \in \check{\sigma}_i(\sigma)$ и константа $\epsilon > 0$, такая что:
\begin{equation}
    u_i(\sigma'_i, \sigma_{-i}) > u_i(\hat{\sigma}_i, \sigma_{-i}) + 3\epsilon
\end{equation}
В силу непрерывности $u_i$ и сходимости $(\sigma^n, \hat{\sigma}^n) \to (\sigma, \hat{\sigma})$, для достаточно большого $n$ справедливо:
\begin{equation}
    u_i(\sigma'_i, \sigma^n_{-i}) > u_i(\sigma'_i, \sigma_{-i}) - \epsilon > u_i(\hat{\sigma}_i, \sigma_{-i}) + 2\epsilon > u_i(\hat{\sigma}^n_i, \sigma^n_{-i}) + \epsilon
\end{equation}
Из полученного неравенства $u_i(\sigma'_i, \sigma^n_{-i}) > u_i(\hat{\sigma}^n_i, \sigma^n_{-i}) + \epsilon$ следует, что $\hat{\sigma}^n_i \notin \check{\sigma}_i(\sigma^n)$. Это прямо противоречит исходному допущению. Таким образом, отображение имеет замкнутый график, и по теореме Какутани неподвижная точка существует.
\end{proof}

\subsection{Свойства равновесия: условие безразличия}

При поиске равновесий в смешанных стратегиях критически важным является следующее свойство, определяющее принцип безразличия.

\begin{lemma}[Условие безразличия]\label{lem:indifference}
Пусть $G = \langle I, S, u \rangle$ — конечная игра, $\sigma^*$ — равновесие Нэша в смешанных стратегиях. Обозначим через $S_i^{\sigma_i^*} \subseteq S_i$ носитель стратегии $\sigma_i^*$ (множество всех чистых стратегий, играемых с положительной вероятностью при смешанной стратегии $\sigma_i^*$). Тогда для всех $s_i \in S_i^{\sigma_i^*}$ выполняется:
\begin{equation}
    u_i(s_i, \sigma^*_{-i}) = u_i(\sigma^*_i, \sigma^*_{-i}),
\end{equation}
а для всех чистых стратегий вне носителя $s'_i \notin S_i^{\sigma_i^*}$:
\begin{equation}
    u_i(s'_i, \sigma^*_{-i}) \le u_i(\sigma^*_i, \sigma^*_{-i}).
\end{equation}
\end{lemma}

Иными словами, в равновесии все чистые стратегии, входящие в состав смешанной стратегии с положительной вероятностью, приносят игроку строго одинаковый ожидаемый выигрыш. Если бы это было не так, игрок мог бы увеличить свой общий ожидаемый выигрыш, перенеся всю вероятность на ту чистую стратегию из носителя, которая дает наибольшую ожидаемую полезность.

\begin{remark}[Инвариантность к аффинным преобразованиям]
Пусть $s^*$ — равновесие в игре $G = \langle I, S, u \rangle$. Рассмотрим игру $G'$, в которой функция полезности $i$-го игрока заменена на $u'_i(s) = \alpha u_i(s) + \beta$, где $\alpha > 0$. Тогда $s^*$ остается равновесием и в игре $G'$. Это правило справедливо как для чистых, так и для смешанных стратегий.
\end{remark}

\subsection{Практические примеры и вычисление равновесий}

\begin{example}[Модель пробития пенальти]
Рассмотрим игру бьющего (Игрок 1) и вратаря (Игрок 2). У каждого есть три чистые стратегии: Лево (L), Центр (C), Право (R). Матрица отражает выигрыши бьющего (вероятность забить гол).

\begin{table}
\centering
\begin{NiceTabular}{cc|ccc}
\toprule
& & \multicolumn{3}{c}{\textbf{Вратарь}} \\
& & \textit{L} & \textit{C} & \textit{R} \\
\midrule
\multirow{3}{*}{\textbf{Бьющий}} 
& \textit{L} & $\alpha$ & $1$ & $1$ \\
& \textit{C} & $\mu$ & $0$ & $\mu$ \\
& \textit{R} & $1$ & $1$ & $\alpha$ \\
\bottomrule
\end{NiceTabular}
\end{table}

Здесь $1 - \alpha$ — вероятность того, что вратарь отразит мяч при угадывании направления (слева или справа), а $1 - \mu$ — вероятность того, что бьющий промажет по центру. Предполагается, что игра антагонистическая (с нулевой суммой).

Анализ симметрии: докажем, что оба игрока выбирают $L$ и $R$ с равными вероятностями. Если допустить, что для бьющего вероятность $q_L > q_R$, то ожидаемая полезность $u_1(L) = 1 - q_L(1 - \alpha) < 1 - q_R(1 - \alpha) = u_1(R)$. Тогда в равновесии бьющий вообще не должен играть $L$ (то есть $p_L = 0$). Но если $p_L = 0$, то для вратаря $u_2(L) < u_2(R)$, что ведет к противоречию $q_L = 0$.
Следовательно, равновесная стратегия имеет вид $(p, 1 - 2p, p)$ либо $(1/2, 0, 1/2)$.

Предположим, вратарь играет стратегию $(q, 1 - 2q, q)$. Если нападающий выбирает $L$, $R$ и $C$ с положительными вероятностями, то по Лемме \ref{lem:indifference}:
\begin{equation}
    u_1(L) = q(1 + \alpha) + 1 - 2q = 2q\mu = u_1(C)
\end{equation}
Отсюда находим равновесную вероятность для вратаря:
\begin{equation}
    q^* = \frac{1}{1 - \alpha + 2\mu}
\end{equation}
Условие $q^* < 1/2$ выполняется при $\mu > \frac{1+\alpha}{2}$. Если $\mu \le \frac{1+\alpha}{2}$, то в равновесии всегда $(p_L, p_C, p_R) = (1/2, 0, 1/2)$.

В итоге, если точность по центру высока ($\mu > \frac{1+\alpha}{2}$):
\begin{equation}
    (p_L, p_C, p_R) = \left( \frac{\mu}{1 - \alpha + 2\mu}, \frac{1 - \alpha}{1 - \alpha + 2\mu}, \frac{\mu}{1 - \alpha + 2\mu} \right)
\end{equation}
\begin{equation}
    (q_L, q_C, q_R) = \left( \frac{1}{1 - \alpha + 2\mu}, \frac{2\mu - 1 - \alpha}{1 - \alpha + 2\mu}, \frac{1}{1 - \alpha + 2\mu} \right)
\end{equation}
\end{example}

На практике интерпретация смешанных стратегий выходит за рамки буквального бросания кубика:
\begin{itemize}
    \item В спорте (теннис, футбол) игроки сознательно действуют непредсказуемо во избежание эксплуатации их паттернов оппонентом.
    \item При большом числе игроков распределение вероятностей может интерпретироваться как доля популяции (например, половина безбилетников идет в первый вагон, половина — во второй).
    \item Согласно концепции очищения Харсаньи, игроки могут играть чистые стратегии на основе приватной информации о малых флуктуациях собственных выигрышей, что для внешнего наблюдателя (или оппонента) выглядит как вероятностное распределение.
\end{itemize}

\begin{example}[Поиск равновесия через огибающие в играх $2 \times N$]
Города $A$ и $B$ соединены тремя дорогами, одна из которых (маршрут $C$) односторонняя. Водитель 1 едет из $A$ в $B$, водитель 2 едет из $B$ в $A$. Время в пути удваивается, если по дороге решают ехать оба водителя (издержки отражены отрицательными выигрышами). 

\begin{table}
\centering
\begin{NiceTabular}{cc|ccc}
\toprule
& & \multicolumn{3}{c}{\textbf{Водитель 2}} \\
& & \textit{T} & \textit{C} & \textit{B} \\
\midrule
\multirow{2}{*}{\textbf{Водитель 1}} 
& \textit{T} & $-8, -8$ & $-4, -5$ & $-4, -6$ \\
& \textit{C} & $-5, -4$ & $-10, -10$ & $-5, -6$ \\
\bottomrule
\end{NiceTabular}
\end{table}

Поиск равновесий в чистых стратегиях дает два исхода: $(T, C)$ и $(C, T)$.
Для поиска равновесий в смешанных стратегиях введем распределение для первого игрока: $\sigma_1 = x \cdot T + (1 - x) \cdot C$, где $x \in [0, 1]$. 
Вычислим ожидаемые полезности для второго игрока как функции от $x$:
\begin{align}
    \mathbb{E}u_2(x, T) &= -8x - 4(1 - x) = -4 - 4x \\
    \mathbb{E}u_2(x, C) &= -5x - 10(1 - x) = -10 + 5x \\
    \mathbb{E}u_2(x, B) &= -6x - 6(1 - x) = -6
\end{align}

Построив верхнюю огибающую данных прямых, мы получим функцию наилучшего ответа второго игрока $\operatorname{BR}_2(x)$:
\begin{equation}
    \operatorname{BR}_2(x) = \begin{cases}
    T, & \text{если } x < 0.5, \\
    \alpha T + (1 - \alpha)B, \; \alpha \in [0, 1], & \text{если } x = 0.5, \\
    B, & \text{если } x \in (0.5, 0.8), \\
    \beta B + (1 - \beta)C, \; \beta \in [0, 1], & \text{если } x = 0.8, \\
    C, & \text{если } x > 0.8.
    \end{cases}
\end{equation}

В равновесии должно выполняться $x \in \operatorname{BR}_1(\operatorname{BR}_2(x))$. Рассмотрим участки смешивания:
\begin{enumerate}
    \item Пусть Игрок 2 смешивает стратегии $T$ и $B$. Это возможно только при $x = 0.5$. Для того чтобы Игрок 1 был готов играть $x = 0.5$ (то есть смешивать $T$ и $C$), он должен быть безразличен между ними. Пусть стратегия Игрока 2 есть $\sigma_2 = \alpha T + (1 - \alpha)B$. Приравниваем выигрыши первого игрока:
    \begin{equation}
        -8\alpha - 4(1 - \alpha) = -5\alpha - 5(1 - \alpha) \iff -4 - 4\alpha = -5 \iff \alpha = \frac{1}{4}
    \end{equation}
    Таким образом, найдено равновесие: $\left(0.5T + 0.5C, \; 0.25T + 0.75B\right)$.
    
    \item Пусть Игрок 2 смешивает стратегии $B$ и $C$. Это происходит при $x = 0.8$. Пусть $\sigma_2 = \beta B + (1 - \beta)C$. Приравниваем выигрыши Игрока 1:
    \begin{equation}
        -4\beta - 4(1 - \beta) = -5\beta - 10(1 - \beta) \iff \beta = 1.2 > 1
    \end{equation}
    Вероятность не может быть больше единицы, следовательно, на этом участке смешанного равновесия не существует.
\end{enumerate}
Итоговое множество равновесий Нэша в игре: $(C, T)$, $(T, C)$ и смешанное $\left(0.5T + 0.5C, \; 0.25T + 0.75B\right)$.
\end{example}
